% =====================================================================
\chapitre{Conception de l'application}
% =====================================================================

Ce chapitre traduit les besoins identifiés en une conception concrète. Il décrit
l'architecture globale de l'application, le modèle de données et ses règles d'accès, les
diagrammes UML des principaux scénarios, le cycle de vie commun à tous les modèles d'IA et,
pour finir, les choix d'interface.

\subsection{Architecture globale}

\equipay repose sur une architecture découplée en trois blocs : l'application mobile, la
plateforme Supabase qui détient les données partagées, et un backend d'IA qui sert les
modèles. Un quatrième bloc, MLflow, ne participe pas au fonctionnement quotidien mais
conserve l'historique des entraînements et désigne la version de chaque modèle que le
backend doit charger.

\subsubsection{Application mobile --- Flutter}

L'application est développée avec Flutter~\cite{flutter} et suit les principes de Material 3.
Elle regroupe les écrans, la logique de présentation et une couche de services (activités,
dépenses, programme, géocodage, IA). Chaque service est défini par une interface, avec une
implémentation réelle (Supabase ou HTTP) et une implémentation en mémoire utilisée par les
tests. Sur iPhone, la carte et la recherche de lieux s'appuient directement sur MapKit~\cite{mapkit}
grâce à une vue native écrite en Swift.

\subsubsection{Données --- Supabase}

Supabase~\cite{supabase} fournit quatre services utilisés par l'application :
l'authentification (e-mail et Google), une base PostgreSQL protégée par des règles RLS, le
stockage des photos et des reçus, et la diffusion en temps réel des modifications. L'application
lit et écrit directement dans la base : aucune donnée ne transite par le backend d'IA.

\subsubsection{Service d'IA --- FastAPI}

Le backend, écrit en Python avec FastAPI~\cite{fastapi}, expose quatre routes, une par
fonction d'IA. Il s'exécute sur un Mac Apple Silicon et fait tourner les modèles avec
MLX~\cite{mlx}. Chaque requête doit porter le jeton JWT de l'utilisateur, que le backend
vérifie avec la clé publique de Supabase. Les modèles sont chargés à la première
utilisation, puis gardés en mémoire.

\subsubsection{Entraînement et suivi --- MLflow}

Les modèles sont entraînés sur un PC équipé d'une RTX 5080, puis importés sur le Mac.
Le serveur MLflow~\cite{mlflow_docs}, installé sur le Mac et accessible sur le réseau local,
reçoit les paramètres, les métriques et les adaptateurs de chaque entraînement. Son registre
attribue l'alias « champion » à la meilleure version, que le backend charge au démarrage.

\subsubsection{Flux global}

La figure~\ref{fig:architecture} montre comment ces blocs communiquent.

\begin{figure}[H]
\centering
\resizebox{\textwidth}{!}{%
\begin{tikzpicture}[
  bloc/.style={rectangle, rounded corners=5pt, draw=#1, line width=1pt, fill=#1!7, align=center,
               text width=4.0cm, inner sep=7pt, font=\small},
  ext/.style={rectangle, rounded corners=5pt, draw=gristexte, dashed, fill=grisfond, align=center,
              text width=3.3cm, inner sep=6pt, font=\small},
  fl/.style={-{Latex[length=2.4mm]}, line width=0.9pt, draw=#1},
  et/.style={font=\scriptsize, fill=white, inner sep=1.5pt, align=center},
  ]
  \node[bloc=lagon] (app) at (0,-2.1) {\textbf{Application mobile}\\[2pt] Flutter (iPhone)\\[4pt]
     \scriptsize Écrans · services\\ \scriptsize carte native MapKit};
  \node[ext] (plans) at (0,-5.4) {\textbf{Apple Plans}\\ \scriptsize recherche de lieux, carte,\\ \scriptsize itinéraires};
  \node[bloc=violet] (sb) at (6.6,1.6) {\textbf{Supabase}\\[2pt] \scriptsize Auth · PostgreSQL + RLS\\
     \scriptsize Storage · Realtime};
  \node[bloc=orange] (api) at (6.6,-2.1) {\textbf{Backend d'IA}\\[2pt] FastAPI sur Mac M3 Pro\\[4pt]
     \scriptsize /chat · /receipts/scan\\ \scriptsize /expenses/describe · /trips/optimize\\[2pt]
     \scriptsize modèles MLX 4 bits};
  \node[ext] (osrm) at (12.6,-1.2) {\textbf{OSRM}\\ \scriptsize temps et distances\\ \scriptsize entre les lieux};
  \node[bloc=recevoir] (mlf) at (12.6,-3.4) {\textbf{MLflow}\\[2pt] \scriptsize essais, métriques,\\ \scriptsize registre « champion »};
  \node[bloc=bleuclair] (pc) at (12.6,-6.2) {\textbf{PC RTX 5080}\\[2pt] \scriptsize entraînement LoRA,\\ \scriptsize évaluation};

  \draw[fl=violet] (app.north) |- node[et, pos=0.75, above=1pt] {données, fichiers, temps réel} (sb.west);
  \draw[fl=orange] (app.east |- api.west) -- node[et, above=2pt] {HTTP + JWT} (api.west);
  \draw[fl=gristexte] (app.south) -- node[et, right=2pt] {MapKit} (plans.north);
  \draw[fl=violet, dashed] (api.north) -- node[et, right=2pt] {clé JWT, historique du chat} (sb.south);
  \draw[fl=gristexte] (api.east |- osrm.west) -- node[et, above=2pt] {temps de trajet} (osrm.west);
  \draw[fl=recevoir] (mlf.west) -- node[et, below=2pt] {version champion} (api.east |- mlf.west);
  \draw[fl=bleuclair] (pc.north) -- node[et, right=2pt] {runs, adaptateurs} (mlf.south);
\end{tikzpicture}}
\caption{Architecture globale et flux de données d'EquiPay}
\label{fig:architecture}
\end{figure}

Le parcours d'une donnée suit toujours la même logique : l'application envoie au backend
une photo, une phrase ou une liste de lieux ; le backend renvoie une proposition structurée
sans rien enregistrer ; l'utilisateur la vérifie, la corrige si besoin, et c'est
l'application qui l'enregistre dans Supabase. Les autres membres du groupe la reçoivent
alors en temps réel.

\begin{encadre}
\textbf{Règle de conception.} L'IA ne crée jamais rien à la place de l'utilisateur : le scanner
et la saisie par phrase pré-remplissent un formulaire, et le trajet proposé n'est enregistré
qu'après un appui sur « Appliquer ». La décision finale reste toujours humaine.
\end{encadre}

% ---------------------------------------------------------------------
\subsection{Conception de la base de données}

Le schéma a été conçu autour de l'activité, à laquelle se rattachent les membres, les
dépenses, les photos et le programme. Il est entièrement décrit par des fichiers de
migration SQL versionnés, appliqués dans l'ordre, ce qui permet de recréer la base à
l'identique. Le tableau~\ref{tab:tables} présente les tables principales.

\begin{table}[H]
\centering
\small
\begin{tabularx}{\textwidth}{|L{4.6cm}|Y|}
\hline
\rowcolor{entete}\textbf{Table} & \textbf{Contenu} \\ \hline
\code{profiles} & Profil de chaque utilisateur : nom d'utilisateur, nom complet, photo. \\ \hline
\code{activities} & Activité : type (voyage, sortie, quotidien, crédit), nom, devise, logo, dates. \\ \hline
\code{activity\_members} & Participants, rôle (créateur, membre, spectateur), date d'archivage personnelle. \\ \hline
\code{activity\_join\_codes}, \code{join\_requests} & Code d'accès à huit caractères et demandes pour rejoindre. \\ \hline
\code{expenses}, \code{expense\_shares} & Dépenses et part de chaque participant, exprimée en pourcentage. \\ \hline
\code{attachments} & Reçus et photos de l'activité ; les fichiers eux-mêmes sont dans le stockage. \\ \hline
\code{trip\_items}, \code{trip\_votes} & Programme du voyage et votes pour ou contre. \\ \hline
\code{chat\_messages} & Historique des échanges avec l'assistant. \\ \hline
\end{tabularx}
\caption{Tables principales de la base de données}
\label{tab:tables}
\end{table}

\subsubsection{Activités, membres et accès}

Chaque activité a ses propres rôles : le créateur gère tout, le membre ajoute des dépenses
et vote, le spectateur consulte. Une même personne peut être créatrice d'une activité et
simple membre d'une autre. L'archivage est personnel : il est stocké dans
\code{activity\_members} et ne masque l'activité que pour celui qui l'archive.

\subsubsection{Dépenses et répartitions}

Une dépense est enregistrée dans sa devise d'origine avec son taux de change vers la devise
de l'activité. La part de chaque participant est stockée en pourcentage, ce qui couvre tous
les modes de répartition (égale, montants, parts) avec une seule représentation. Une
fonction SQL vérifie que les parts d'une dépense totalisent exactement 100 \%, et une
fonction \code{save\_expense} enregistre la dépense et ses parts en une seule transaction.

\subsubsection{Programme du voyage}

La table \code{trip\_items} contient les lieux proposés, avec leur adresse, leur date,
leur heure et leurs coordonnées, remplies automatiquement à partir de l'adresse. Les votes
sont stockés dans \code{trip\_votes}, avec une seule ligne par personne et par lieu. Une
fonction \code{has\_program} réserve le programme aux activités de type voyage ou sortie.

\subsubsection{Sécurité au niveau des lignes}

Toutes les tables sont protégées par des règles RLS. Elles reposent sur quelques fonctions
réutilisées partout (\code{is\_member}, \code{is\_owner}, \code{can\_edit}) : un
utilisateur ne lit que les activités dont il est membre, ne modifie que ce qu'il a créé
(sauf le créateur de l'activité), et un spectateur ne peut rien modifier. Ces règles
s'appliquent quelle que soit la façon dont la base est interrogée.

% ---------------------------------------------------------------------
\subsection{Diagrammes UML}

Cette section présente les diagrammes UML qui modélisent l'application : le diagramme de
cas d'utilisation décrit les interactions entre les acteurs et le système, le diagramme de
classes représente la structure des données, et les diagrammes de séquence détaillent le
déroulement des quatre fonctions d'IA.

\subsubsection{Diagramme de cas d'utilisation}

\begin{figure}[H]
\centering
\resizebox{0.98\textwidth}{!}{%
\begin{tikzpicture}[
  uc/.style={ellipse, draw=lagon!80!black, fill=lagonclair!40, line width=0.7pt, minimum width=3.9cm,
             minimum height=0.8cm, font=\scriptsize, align=center, inner sep=1pt},
  ucx/.style={uc, draw=orange!80!black, fill=orangeclair!60},
  as/.style={line width=0.6pt, draw=black!65},
  dep/.style={dashed, -{Latex[length=2mm]}, draw=black!60, line width=0.6pt},
  st/.style={font=\tiny\itshape, fill=white, inner sep=1pt},
  ]
  % Acteur : bonhomme
  \newcommand{\acteur}[3]{%
    \begin{scope}[shift={(#1)}]
      \draw[line width=0.8pt] (0,0.55) circle (0.2);
      \draw[line width=0.8pt] (0,0.35) -- (0,-0.25) (-0.32,0.15) -- (0.32,0.15) (0,-0.25) -- (-0.25,-0.65) (0,-0.25) -- (0.25,-0.65);
      \node[font=\small\bfseries, align=center, anchor=north] at (0,-0.75) {#2};
      \coordinate (#3) at (0,0);
    \end{scope}}
  % Limite du système
  \draw[rounded corners=6pt, line width=0.9pt, draw=bleu, fill=grisfond!50] (1.25,0.35) rectangle (10.75,-10.55);
  \node[font=\bfseries\color{bleu}, anchor=north west] at (1.4,0.3) {EquiPay};

  % Cas du membre (colonne de gauche)
  \foreach \n/\t [count=\i] in {
      auth/{S'inscrire, se connecter},
      creer/{Créer une activité},
      rejoindre/{Rejoindre par code},
      depense/{Ajouter une dépense},
      soldes/{Consulter soldes et résumé},
      lieu/{Ajouter un lieu au programme},
      voter/{Voter pour ou contre},
      trajet/{Organiser le trajet},
      question/{Poser une question à l'assistant}}
    \node[uc] (\n) at (3.6,-0.9*\i-0.1) {\t};
  \node[uc] (accepter) at (3.6,-9.95) {Accepter une demande};

  % Extensions et inclusions (colonne de droite)
  \node[ucx] (scan) at (8.4,-2.75) {Scanner un ticket};
  \node[ucx] (phrase) at (8.4,-3.85) {Décrire la dépense\\ en une phrase};
  \node[ucx] (carte) at (8.4,-5.5) {Placer le lieu sur la carte};
  \node[ucx] (appliquer) at (8.4,-8.95) {Appliquer le programme};

  \draw[dep] (scan) -- node[st, above, sloped] {«extend»} (depense);
  \draw[dep] (phrase) -- node[st, below, sloped] {«extend»} (depense);
  \draw[dep] (lieu) -- node[st, above] {«include»} (carte);
  \draw[dep] (appliquer) -- node[st, below, sloped] {«extend»} (trajet);

  \acteur{-0.9,-3.9}{Membre}{membre}
  \acteur{-0.9,-8.6}{Créateur}{createur}
  \acteur{12.9,-5.0}{Modèles\\ d'IA}{ia}

  \draw[-{Triangle[open, length=3mm]}, line width=0.7pt] ([yshift=0.9cm]createur) -- ([yshift=-1.75cm]membre);
  \foreach \n in {auth, creer, rejoindre, depense, soldes, lieu, voter, trajet, question}
    \draw[as] ([xshift=0.35cm]membre) -- (\n.west);
  \draw[as] ([xshift=0.35cm]createur) -- (accepter.west);
  \draw[as] ([xshift=0.35cm, yshift=-0.1cm]createur) -- (appliquer.west);
  \foreach \n in {scan, phrase} \draw[as] ([xshift=-0.35cm]ia) -- (\n.east);
  \draw[as] ([xshift=-0.35cm]ia) -- (trajet.east);
  \draw[as] ([xshift=-0.35cm]ia) -- (question.east);
\end{tikzpicture}}
\caption{Diagramme de cas d'utilisation de l'application}
\label{fig:cas_utilisation}
\end{figure}

La figure~\ref{fig:cas_utilisation} présente les interactions entre les acteurs et le
système. Le membre gère son compte, crée ou rejoint des activités, ajoute des dépenses,
consulte les soldes, propose des lieux, vote, demande l'organisation du trajet et
interroge l'assistant. Le créateur est un membre qui dispose en plus de deux droits :
accepter les demandes pour rejoindre et appliquer le programme proposé. Les modèles d'IA
interviennent comme acteur secondaire ; la carte et le calcul des trajets s'appuient en
plus sur Apple Plans et OSRM, qui ne sont pas représentés pour alléger le schéma.

\subsubsection{Diagramme de classes}

\begin{figure}[H]
\centering
\resizebox{\textwidth}{!}{%
\begin{tikzpicture}[
  classe/.style={rectangle split, rectangle split parts=2, draw=bleu, line width=0.7pt,
                 rectangle split part fill={bleu!12, white}, text width=3.55cm, align=left,
                 font=\scriptsize, inner sep=3.5pt},
  as/.style={line width=0.7pt, draw=bleu!80},
  comp/.style={as, {Diamond[length=3.2mm, width=2.2mm]}-},
  mult/.style={font=\tiny, inner sep=1pt},
  ]
  \newcommand{\titre}[1]{\centering\textbf{\small #1}}
  \node[classe] (jr) at (0,0) {\titre{JoinRequest}\nodepart{two}
     id : uuid\\ activity\_name : text\\ status : pending | accepted\\ \hspace{2.6em}| declined\\ created\_at, decided\_at};
  \node[classe] (pr) at (5.4,0) {\titre{Profile}\nodepart{two}
     id : uuid\\ username : text\\ email : text\\ full\_name : text\\ avatar\_url : text};
  \node[classe] (ch) at (10.8,0) {\titre{ChatMessage}\nodepart{two}
     id : uuid\\ role : user | assistant\\ content : text\\ created\_at : timestamptz};

  \node[classe] (ex) at (0,-4.7) {\titre{Expense}\nodepart{two}
     id : uuid\\ name : text\\ amount : numeric\\ currency : char(3)\\ exchange\_rate : numeric\\ expense\_date : date\\ paid\_by : Profile};
  \node[classe] (ac) at (5.4,-4.7) {\titre{Activity}\nodepart{two}
     id : uuid\\ type : voyage | sortie |\\ \hspace{2.6em}quotidien | crédit\\ name : text\\ default\_currency : char(3)\\ start\_date, end\_date\\ cover\_url : text};
  \node[classe] (am) at (10.8,-4.7) {\titre{ActivityMember}\nodepart{two}
     role : owner | member\\ \hspace{2.6em}| viewer\\ joined\_at : timestamptz\\ archived\_at : timestamptz};

  \node[classe] (es) at (-0.6,-9.6) {\titre{ExpenseShare}\nodepart{two}
     user\_id : Profile\\ percentage : numeric};
  \node[classe] (at) at (3.4,-9.6) {\titre{Attachment}\nodepart{two}
     id : uuid\\ kind : receipt | photo\\ storage\_path : text};
  \node[classe] (tv) at (7.4,-9.6) {\titre{TripVote}\nodepart{two}
     user\_id : Profile\\ vote : -1 | 1};
  \node[classe] (ti) at (11.4,-9.6) {\titre{TripItem}\nodepart{two}
     id : uuid\\ name, place, address\\ planned\_on : date\\ planned\_at : time\\ latitude, longitude\\ done : boolean};

  \draw[as] (jr) -- node[mult, pos=0.12, above] {*} node[mult, pos=0.88, above] {1} (pr);
  \draw[as] (pr) -- node[mult, pos=0.12, above] {1} node[mult, pos=0.88, above] {*} (ch);
  \draw[as] (pr) -- node[mult, pos=0.15, right] {1} node[mult, pos=0.5, right, font=\tiny\itshape] {crée} node[mult, pos=0.85, right] {*} (ac);
  \draw[as] (jr.south east) -- node[mult, pos=0.1, right] {*} node[mult, pos=0.9, left] {1} (ac.north west);
  \draw[as] (pr.south east) -- node[mult, pos=0.1, right] {1} node[mult, pos=0.9, above] {*} (am.north west);
  \draw[comp] (ac) -- node[mult, pos=0.85, above] {*} (am);
  \draw[comp] (ac) -- node[mult, pos=0.85, above] {*} (ex);
  \draw[comp] (ex) -- node[mult, pos=0.85, left] {1..*} (es.north -| ex);
  \draw[comp] (ac.south west) -- node[mult, pos=0.85, right] {*} (at.north);
  \draw[as] (ex.south east) -- node[mult, pos=0.1, above] {0..1} node[mult, pos=0.85, left] {*} (at.north west);
  \draw[comp] (ac.south east) -- node[mult, pos=0.88, above] {*} (ti.north west);
  \draw[comp] (ti) -- node[mult, pos=0.8, above] {*} (tv);
\end{tikzpicture}}
\caption{Diagramme de classes du modèle de données}
\label{fig:classes}
\end{figure}

La figure~\ref{fig:classes} modélise la structure des données, directement dérivée du
schéma PostgreSQL. L'activité est au centre : elle se compose de ses membres, de ses
dépenses, de ses photos et de son programme, qui disparaissent avec elle. Chaque dépense
se décompose en parts, une par participant, et peut être liée à la photo de son reçu.
Les demandes pour rejoindre et l'historique du chat sont rattachés au profil de
l'utilisateur.

\subsubsection{Diagrammes de séquence}

Les diagrammes suivants détaillent le déroulement des quatre fonctions d'IA, de l'action
de l'utilisateur jusqu'à l'enregistrement du résultat.

\begin{figure}[H]
\centering
\resizebox{0.97\textwidth}{!}{\footnotesize%
\begin{sequencediagram}
  \gdef\unitfactor{0.82}
  \tikzstyle{inststyle}+=[fill=lagonclair!50, rounded corners=2pt]
  \newthread[lagonclair]{u}{:Utilisateur}
  \newinst[2.9]{app}{:Application}
  \newinst[3.3]{api}{:Backend FastAPI}
  \newinst[3.3]{vlm}{:Qwen3-VL-4B (MLX)}
  \newinst[2.5]{sb}{:Supabase}
  \begin{call}{u}{photographie le ticket}{app}{formulaire pré-rempli}
    \begin{call}{app}{POST /receipts/scan}{api}{champs de la dépense}
      \begin{callself}{api}{vérifie le jeton}{}
      \end{callself}
      \begin{call}{api}{lit l'image}{vlm}{JSON brut}
      \end{call}
      \begin{callself}{api}{normalise montant, devise, date}{}
      \end{callself}
    \end{call}
  \end{call}
  \begin{call}{u}{vérifie, corrige et valide}{app}{dépense ajoutée}
    \begin{call}{app}{save\_expense}{sb}{ok}
    \end{call}
  \end{call}
  \mess[1]{sb}{temps réel : nouvelle dépense}{app}
\end{sequencediagram}}
\caption{Diagramme de séquence --- Scanner un ticket de caisse}
\label{fig:seq_scan}
\end{figure}

La figure~\ref{fig:seq_scan} illustre la lecture d'un ticket : l'application envoie la photo
au backend, qui vérifie le jeton de l'utilisateur, fait lire l'image par le modèle
vision-langage puis normalise sa réponse. Le formulaire est pré-rempli, et la dépense n'est
enregistrée dans Supabase qu'après validation par l'utilisateur.

\begin{figure}[H]
\centering
\resizebox{0.97\textwidth}{!}{\footnotesize%
\begin{sequencediagram}
  \gdef\unitfactor{0.82}
  \tikzstyle{inststyle}+=[fill=lagonclair!50, rounded corners=2pt]
  \newthread[lagonclair]{u}{:Utilisateur}
  \newinst[2.9]{app}{:Application}
  \newinst[3.3]{api}{:Backend FastAPI}
  \newinst[3.3]{llm}{:equipay-describe (MLX)}
  \newinst[2.5]{sb}{:Supabase}
  \begin{call}{u}{écrit sa phrase}{app}{formulaire pré-rempli}
    \begin{call}{app}{POST /expenses/describe}{api}{dépense structurée}
      \begin{call}{api}{phrase, membres, date}{llm}{JSON brut}
      \end{call}
      \begin{callself}{api}{relie les noms aux membres}{}
      \end{callself}
      \begin{callself}{api}{parts ramenées à 100 \%}{}
      \end{callself}
    \end{call}
  \end{call}
  \begin{call}{u}{vérifie et valide}{app}{dépense ajoutée}
    \begin{call}{app}{save\_expense}{sb}{ok}
    \end{call}
  \end{call}
\end{sequencediagram}}
\caption{Diagramme de séquence --- Décrire une dépense en une phrase}
\label{fig:seq_phrase}
\end{figure}

La figure~\ref{fig:seq_phrase} décrit la saisie par phrase. Le modèle reçoit la phrase, la
liste des membres et la date du jour ; le backend relie ensuite chaque nom cité à un membre
réel, même mal orthographié, et convertit la répartition en pourcentages dont la somme fait
exactement 100.

\begin{figure}[H]
\centering
\resizebox{0.97\textwidth}{!}{\footnotesize%
\begin{sequencediagram}
  \gdef\unitfactor{0.82}
  \tikzstyle{inststyle}+=[fill=lagonclair!50, rounded corners=2pt]
  \newthread[lagonclair]{u}{:Utilisateur}
  \newinst[2.7]{app}{:Application}
  \newinst[3.1]{api}{:Backend FastAPI}
  \newinst[2.7]{osrm}{:OSRM}
  \newinst[2.7]{llm}{:Qwen3-4B (MLX)}
  \newinst[2.0]{sb}{:Supabase}
  \begin{call}{u}{choisit les options}{app}{trajet proposé}
    \begin{call}{app}{POST /trips/optimize}{api}{ordre, heures, gain, texte}
      \begin{call}{api}{matrice des temps entre lieux}{osrm}{durées, distances}
      \end{call}
      \begin{callself}{api}{répartit par jour, meilleur ordre}{}
      \end{callself}
      \begin{call}{api}{programme calculé à rédiger}{llm}{texte}
      \end{call}
    \end{call}
  \end{call}
  \begin{call}{u}{appliquer (créateur)}{app}{programme mis à jour}
    \begin{call}{app}{date et heure de chaque lieu}{sb}{ok}
    \end{call}
  \end{call}
\end{sequencediagram}}
\caption{Diagramme de séquence --- Organiser le trajet d'une journée}
\label{fig:seq_trajet}
\end{figure}

La figure~\ref{fig:seq_trajet} présente l'organisation du trajet. Les temps de parcours sont
calculés par OSRM et l'ordre de visite par un algorithme ; le modèle de langage ne fait que
rédiger le programme à partir de ces résultats, sans pouvoir modifier l'ordre ni les heures.
Seul le créateur de l'activité peut appliquer le programme proposé.

\begin{figure}[H]
\centering
\resizebox{0.86\textwidth}{!}{\footnotesize%
\begin{sequencediagram}
  \gdef\unitfactor{0.82}
  \tikzstyle{inststyle}+=[fill=lagonclair!50, rounded corners=2pt]
  \newthread[lagonclair]{u}{:Utilisateur}
  \newinst[2.9]{app}{:Application}
  \newinst[3.3]{api}{:Backend FastAPI}
  \newinst[3.3]{llm}{:equipay-assistant (MLX)}
  \newinst[2.5]{sb}{:Supabase}
  \begin{call}{u}{pose une question}{app}{réponse affichée}
    \begin{call}{app}{POST /chat (question, JWT)}{api}{réponse}
      \begin{call}{api}{question + prompt système}{llm}{réponse générée}
      \end{call}
      \begin{call}{api}{enregistre l'échange}{sb}{ok}
      \end{call}
    \end{call}
  \end{call}
\end{sequencediagram}}
\caption{Diagramme de séquence --- Question à l'assistant}
\label{fig:seq_chat}
\end{figure}

La figure~\ref{fig:seq_chat} montre l'échange avec l'assistant. Le backend génère la réponse
avec le modèle champion puis enregistre la question et la réponse dans Supabase, au nom de
l'utilisateur, afin de conserver l'historique de la conversation.

% ---------------------------------------------------------------------
\subsection{Architecture fonctionnelle des modules}

L'application s'articule autour de quatre modules fonctionnels, chacun responsable d'une
partie du parcours de l'utilisateur.

\subsubsection{Module des dépenses et des soldes}

Ce module gère les activités, les dépenses et leur répartition. Le solde de chaque
participant est la différence entre ce qu'il a payé et la somme de ses parts, convertie
dans la devise de l'activité. À partir de ces soldes, l'application calcule la liste des
remboursements qui remettent tous les comptes à zéro avec le moins de virements possible,
et propose des graphiques et des statistiques sur les dépenses du groupe.

\subsubsection{Module de saisie assistée}

Ce module regroupe les deux façons d'ajouter une dépense sans tout taper : le scanner de
tickets et la description en une phrase. Les deux aboutissent au même formulaire, où chaque
champ rempli par l'IA est signalé avec son origine (« Lu sur le ticket », « Compris dans ta
phrase ») et où les champs manquants sont mis en évidence.

\subsubsection{Module du programme et du trajet}

Ce module couvre la liste des lieux à visiter, les votes, l'affichage en liste, en
calendrier et sur une carte, ainsi que l'organisation du trajet. Les lieux sont placés
automatiquement sur la carte à partir de leur nom ou de leur adresse grâce à la recherche
d'Apple Plans.

\subsubsection{Module d'assistance}

L'assistant est accessible depuis une bulle « Support » présente sur les écrans principaux.
Il répond aux questions sur le fonctionnement de l'application, oriente vers le support
quand il ne sait pas, et refuse poliment les questions sans rapport avec \equipay.

% ---------------------------------------------------------------------
\subsection{Conception de la chaîne d'entraînement et de mise en service}

Les trois modèles affinés suivent le même cycle de vie, représenté dans la
figure~\ref{fig:mlops}. Ce cycle garantit qu'une nouvelle version n'est servie à
l'application que si elle fait mieux que la précédente sur le même jeu de test.

\begin{figure}[H]
\centering
\begin{tikzpicture}[
  etape/.style={rectangle, rounded corners=6pt, fill=#1, text=white, align=center, font=\small\bfseries,
                text width=3.3cm, minimum height=1.35cm, drop shadow={opacity=0.18}},
  det/.style={font=\scriptsize, align=center, text width=3.5cm, text=black!75},
  fl/.style={-{Latex[length=2.6mm]}, line width=1.1pt, draw=gristexte},
  ]
  \node[etape=lagon]     (d) at (0,2.2)     {1. Données\\ versionnées};
  \node[etape=violet]    (t) at (5.1,2.2)   {2. Entraînement\\ LoRA sur GPU};
  \node[etape=orange]    (e) at (10.2,2.2)  {3. Évaluation\\ comparative};
  \node[etape=recevoir]  (r) at (10.2,-1.6) {4. Registre MLflow\\ alias « champion »};
  \node[etape=bleuclair] (c) at (5.1,-1.6)  {5. Conversion\\ MLX 4 bits};
  \node[etape=ambre]     (s) at (0,-1.6)    {6. Service\\ FastAPI};
  \node[det, below=2pt of d] {génération, empreinte\\ du jeu enregistrée};
  \node[det, below=2pt of t] {PC RTX 5080, arrêt\\ anticipé sur la validation};
  \node[det, below=2pt of e] {base, ancienne et nouvelle\\ version sur le même test};
  \node[det, below=2pt of r] {seulement si la nouvelle\\ version est meilleure};
  \node[det, below=2pt of c] {fusion de l'adaptateur,\\ réévaluation sur le Mac};
  \node[det, below=2pt of s] {chargement au premier\\ appel, PyTorch en secours};
  \draw[fl] (d) -- (t);
  \draw[fl] (t) -- (e);
  \draw[fl] (e) -- (r);
  \draw[fl] (r) -- (c);
  \draw[fl] (c) -- (s);
  \draw[fl, dashed] (s.north) -- node[font=\scriptsize\itshape, fill=white, inner sep=2pt, align=center] {limites\\ observées} (d.south);
\end{tikzpicture}
\caption{Cycle de vie d'un modèle, des données à la mise en service}
\label{fig:mlops}
\end{figure}

\subsubsection{Construction des jeux de données}

Aucun jeu public ne correspondait à nos besoins pour l'assistant et la saisie par phrase.
Nous les avons donc construits à partir de sources écrites à la main (questions, réponses
de référence, phrases types) puis enrichis par génération de variantes. Pour le scanner, les
jeux publics de tickets ont été convertis dans notre format et complétés par des tickets
synthétiques dans les devises qui nous intéressent. Chaque jeu est découpé en
entraînement, validation et test, et son empreinte est enregistrée avec chaque essai.

\subsubsection{Entraînement}

Les modèles sont affinés par LoRA sur le PC, avec la bibliothèque PEFT~\cite{peft}. La
perte sur le jeu de validation est mesurée régulièrement, et l'entraînement s'arrête
lorsqu'elle cesse de s'améliorer. C'est l'adaptateur du meilleur point de validation qui est
conservé, et non le dernier.

\subsubsection{Évaluation et choix du champion}

Chaque nouvelle version est comparée au modèle de base et à la version précédente, sur le
même jeu de test et avec exactement le même code de calcul des métriques. Les métriques sont
regroupées en un score global pondéré, propre à chaque tâche, qui sert à décider. Si la
nouvelle version l'emporte, elle reçoit l'alias « champion » dans le registre MLflow.

\subsubsection{Conversion et mise en service}

L'adaptateur est fusionné avec le modèle de base, puis le modèle obtenu est converti en
4 bits avec MLX et réévalué sur le Mac, pour mesurer précisément ce que coûte la
quantification. Le backend charge ensuite la version champion. Pour le scanner, la version
PyTorch en pleine précision est conservée et sert de secours automatique si le chargement
MLX échoue.

% ---------------------------------------------------------------------
\subsection{Conception des interfaces utilisateur}

\subsubsection{Charte graphique}

L'identité visuelle, baptisée « Lagon », associe un sarcelle apaisant, qui évoque la
confiance et le voyage, à un orange énergique réservé aux actions principales. Les couleurs
des soldes sont les mêmes partout : vert pour ce que l'on doit recevoir, rouge pour ce que
l'on doit rembourser, gris lorsque l'on est à jour. Chaque type d'activité a sa propre
couleur (figure~\ref{fig:charte}). Le logo, un « e » sarcelle qui s'envole en avion au-dessus
d'un porte-monnaie orange, réunit le voyage et le partage des dépenses.

\begin{figure}[H]
\centering
\begin{tikzpicture}[
  pastille/.style={rectangle, rounded corners=4pt, minimum width=1.75cm, minimum height=1.05cm, fill=#1},
  nom/.style={font=\scriptsize\bfseries, align=center, text width=1.9cm},
  hex/.style={font=\tiny\ttfamily, text=gristexte},
  ]
  \foreach \c/\n/\h [count=\i] in {lagon/Principale/0D9488, orange/Accent/F97316, recevoir/À recevoir/16A34A,
                                  rembourser/À rembourser/DC2626, gristexte/À jour/64748B}
  {
    \node[pastille=\c] (p\i) at (2.15*\i,0) {};
    \node[nom, below=3pt of p\i] (n\i) {\n};
    \node[hex, below=-1pt of n\i] {\#\h};
  }
  \foreach \c/\n/\h [count=\i] in {lagon/Voyage/0D9488, violet/Sortie/8B5CF6, bleuclair/Quotidien/3B82F6, ambre/Crédit/F59E0B}
  {
    \node[pastille=\c, minimum width=1.5cm] (q\i) at (1.2+2.6*\i,-2.65) {};
    \node[nom, below=3pt of q\i] (m\i) {\n};
    \node[hex, below=-1pt of m\i] {\#\h};
  }
  \node[font=\small\bfseries\color{bleu}, anchor=east] at (1.0,0) {Couleurs};
  \node[font=\small\bfseries\color{bleu}, anchor=east, align=right] at (1.0,-2.65) {Types\\ d'activité};
\end{tikzpicture}
\caption{Palette « Lagon » : couleurs principales et couleurs des types d'activité}
\label{fig:charte}
\end{figure}

\subsubsection{Organisation des écrans}

La navigation repose sur une barre d'onglets en bas de l'écran : les activités, les demandes
pour rejoindre et le profil (un quatrième onglet est réservé à une évolution future). Depuis une activité, l'utilisateur accède à ses dépenses, à ses
membres, à ses photos et à son programme ; le menu à trois points regroupe les actions moins
fréquentes (modifier, archiver, résumé, graphiques, statistiques, supprimer). La
figure~\ref{fig:navigation} résume ce parcours.

\begin{figure}[H]
\centering
\resizebox{\textwidth}{!}{%
\begin{tikzpicture}[
  ecr/.style={rectangle, rounded corners=4pt, draw=#1, line width=0.8pt, fill=#1!8, font=\scriptsize,
              align=center, minimum height=0.75cm, text width=2.5cm},
  fl/.style={-{Latex[length=2mm]}, draw=gristexte, line width=0.7pt},
  ]
  \node[ecr=bleu] (tab) at (0,0) {\textbf{Barre d'onglets}};
  \node[ecr=lagon] (acc) at (-4.4,-1.5) {Accueil\\ (mes activités)};
  \node[ecr=lagon] (dem) at (0,-1.5) {Demandes\\ reçues / envoyées};
  \node[ecr=lagon] (pro) at (4.4,-1.5) {Profil\\ et préférences};
  \node[ecr=violet] (act) at (-4.4,-3.2) {Page de l'activité};
  \node[ecr=orange] (dep) at (-8.2,-4.9) {Dépenses\\ (ajout, détail)};
  \node[ecr=orange] (mem) at (-5.4,-4.9) {Membres et photos};
  \node[ecr=orange] (prg) at (-2.6,-4.9) {Programme};
  \node[ecr=orange] (men) at (0.4,-4.9) {Menu à trois points\\ résumé, graphiques};
  \node[ecr=ambre] (aj) at (-8.2,-6.6) {Saisie, scanner,\\ phrase};
  \node[ecr=ambre] (vues) at (-4.0,-6.6) {Liste, calendrier,\\ carte};
  \node[ecr=ambre] (traj) at (-1.2,-6.6) {Trajet proposé};
  \node[ecr=recevoir] (sup) at (4.4,-3.2) {Assistant\\ (bulle Support)};
  \foreach \a/\b in {tab/acc, tab/dem, tab/pro, acc/act, act/dep, act/mem, act/prg, act/men, dep/aj, prg/vues, prg/traj, pro/sup}
    \draw[fl] (\a) -- (\b);
\end{tikzpicture}}
\caption{Organisation des principaux écrans de l'application}
\label{fig:navigation}
\end{figure}

\subsubsection{Expérience utilisateur}

L'expérience a été pensée pour un usage rapide, souvent debout ou en déplacement. Les
actions principales sont toujours au même endroit (bouton orange en bas à droite), chaque
action donne un retour immédiat (chargement, confirmation, erreur expliquée), et les écrans
vides indiquent ce qu'il faut faire pour commencer. L'application propose un mode sombre,
fonctionne hors ligne en affichant un bandeau, et se met à jour d'elle-même lorsqu'un autre
membre ajoute une dépense ou vote.

% ---------------------------------------------------------------------
\subsection{Conclusion du chapitre}

Ce chapitre a posé les bases de la conception d'\equipay : une architecture qui sépare
clairement l'application, les données et l'IA, un modèle de données protégé par des règles
d'accès strictes, des diagrammes UML qui formalisent les scénarios clés, et un cycle de vie
commun à tous les modèles. Cette conception sert de fondation à la réalisation présentée
dans le chapitre suivant.
